\documentclass{beamer}
\usepackage{moreverb} 
\usepackage{listings}
\usepackage{mflogo}
% imprimir
% \documentclass[handout]{beamer} 
% \usepackage{pgfpages}
% \pgfpagesuselayout{4 on 1}[a4paper,landscape,border shrink=5mm]

\mode<presentation> {
  \usetheme{Warsaw}
  \setbeamercovered{opaque}
}

\usebackgroundtemplate{\includegraphics[width=\paperwidth]{format/libresoft-bg.png}}
\usepackage[spanish]{babel}
\usepackage[utf8]{inputenc}
\usepackage{graphics}
\usepackage{amssymb} % Simbolos matematicos

%% Metadatos del PDF.
\hypersetup{  
  pdftitle={Almacenamiento: ejercicios},
  pdfauthor={Miguel Vidal},
  pdfcreator={GSyC/Libresoft},
  pdfproducer=PDFLaTeX,
  pdfsubject={Master on Free Software},
}
%%

\defbeamertemplate*{footline}{shadow theme}
{%
  \leavevmode%
  \hbox{\begin{beamercolorbox}[wd=.5\paperwidth,ht=2.5ex,dp=1.125ex,leftskip=.1cm plus1fil,rightskip=.2cm]{author in head/foot}%
\includegraphics[scale=0.40]{format/cc-by-80x15.png} \hspace{0.05cm}
	Miguel Vidal / Jose Castro
  \end{beamercolorbox}%
  \begin{beamercolorbox}[wd=.5\paperwidth,ht=2.5ex,dp=1.125ex,leftskip=.3cm,rightskip=.3cm plus1fil]{title in head/foot}%
    \usebeamerfont{title in head/foot}\insertshorttitle%
    \hspace{2cm} \usebeamerfont{author in head/foot}\insertframenumber\,/\,\inserttotalframenumber\hfill
  \end{beamercolorbox}}%
  \vskip0pt%
}


\begin{document}

\title{Almacenamiento: ejercicios}
\subtitle{System Integration}
% \institute{\texttt{http://gsyc.urjc.es/\~{}mvidal} \\ Twitter: \texttt{@mvidallopez}}
\author{Miguel Vidal, Jose Castro} 
\date{\footnotesize{Master on Free Software \\ March 30th, 2012}}
\author{Miguel Vidal \hspace{1cm} Jose Castro \\
\hspace{0.5mm} {\tiny Twitter: @mvidallopez \hspace{1.1cm}Twitter: @jfcastroluis}
}


\frame{
\maketitle
\begin{center}
\includegraphics[width=6cm]{format/gsyc-urjc}
\end{center}
}

%% License slide
\begin{frame}
  \vspace{2cm}
  \begin{flushright}
    {\small \copyright{} 2010-2012 Miguel Vidal, Jose Castro} \\
%    \vspace{0.25cm}
    \medskip
    {\scriptsize This work is licensed under \\ a Creative Commons Attribution 3.0 License}
%    \vspace{0.10cm}
  \end{flushright}
  \begin{flushright}
    \href{http://creativecommons.org/licenses/by/3.0/es}{\includegraphics[width=2cm]{format/cc-by.png}} \\
    {\tiny \url{http://creativecommons.org/licenses/by/3.0}}
  \end{flushright}
\end{frame}%%

\usebackgroundtemplate{}

\AtBeginSection[]
{
\begin{frame}<beamer>
\begin{center}
{\huge \insertsection}
\end{center}
\end{frame}
}


\AtBeginSubsection[]
{
  \begin{frame}<beamer>{Índice}
    \tableofcontents[currentsection,currentsubsection]
  \end{frame}
}

%%

\normalsize

% Inicio de la presentacion

\section{VirtualBox}
\begin{frame}
  \frametitle{VirtualBox}

  \begin{block}{¿Qué es una VM?}
Un proceso en el host que maneja un fichero con una configuración de dispositivos emulados: memoria, red, sonido...
  \end{block}

  \begin{itemize}
    \item Los discos en Virtualbox tienen formato VDI y están asociados a una VM. No son más que un fichero del host.
    \item Los discos pueden ser \alert{dynamically allocated} (asignados dinámicamente) o \alert{Fixed size} (imagen de tamaño fijo)
  \end{itemize}
\end{frame}

\begin{frame}
  \frametitle{VirtualBox: configuración de red de una VM}

Los tipos de red en Virtualbox:

  \begin{itemize}
    \item \alert{NAT:} la máquina sale hacia afuera pero no es accesible ni desde fuera ni desde el host.
    \item \alert{Bridge:} se asocia con un bridge a una interfaz y es una máquina más de la red.
    \item \alert{Internal:} se crea una red virtual entre las máquinas de esa red.
    \item \alert{Host-only:} las máquinas tienen conexión hacia afuera y desde el host.
     En este caso, podemos configurar algunas cosas desde las preferencias
  \end{itemize}
\end{frame}

%%%%%%%%%%%%%%%%%%%%%%%%%%%%%%%%%%%%%%%%%%%%%%%%%%%%%%%%%%%%%%%%%%%%%%%%%

\section{Ejercicios}

%%%%%%%%%%%%%%%%%%%%%%%%%%%%%%%%%%%%%%%%%%%%%%%%%%%%%%%%%%%%%%%%%%%%%%%%%

\subsection{RAIDs}

\begin{frame}
  \frametitle{RAIDs: Ejercicio 1}

\begin{itemize}
\item \textbf{Creación de un RAID-1 con 2 dispositivos de bloques creados en VirtualBox.}
\end{itemize}

\end{frame}

%%%%%%%%%%%%%%%%%%%%%%%%%%%%%%%%%%%%%%%%%%%%%%%%%%%%%%%%%%%%%%%%%%%%%%%%%

\begin{frame}
  \frametitle{RAIDs: Ejercicio 1}

\begin{itemize}
\item \textbf{Creación de un RAID-1 con 2 dispositivos de bloques creados en VirtualBox.}
\end{itemize}

\pause

  \begin{block}{}
\# apt-get install mdadm
  \end{block}
Seleccionamos \texttt{none} y que lo arranque automáticamente.

  \begin{block}{}
\# cat /proc/mdstat \\

\# ls /dev/sd*

\# mdadm --create /dev/md0 --level=raid1 --raid-devices=2 $\backslash$ \\ \hspace{1cm} /dev/sda /dev/sdb \\
\# cat /proc/mdstat \\
\# mdadm --detail /dev/md0
  \end{block}
\end{frame}


%%%%%%%%%%%%%%%%%%%%%%%%%%%%%%%%%%%%%%%%%%%%%%%%%%%%%%%%%%%%%%%%%%%%%%%%%

\begin{frame}
  \frametitle{RAIDs: Ejercicio 2}

\begin{enumerate}
\item \textbf{Formatear el dispositivo RAID con el sistema de ficheros favorito y montarlo en \texttt{/mnt}}
\item \textbf{Crear un fichero de texto de nombre \texttt{foo} y contenido \texttt{bar}}
\end{enumerate}

\end{frame}

%%%%%%%%%%%%%%%%%%%%%%%%%%%%%%%%%%%%%%%%%%%%%%%%%%%%%%%%%%%%%%%%%%%%%%%%%

\begin{frame}
  \frametitle{RAIDs: Ejercicio 2}

\begin{enumerate}
\item \textbf{Formatear el dispositivo RAID con el sistema de ficheros favorito y montarlo en \texttt{/mnt}}
\item \textbf{Crear un fichero de texto de nombre \texttt{foo} y contenido \texttt{bar}}
\end{enumerate}

\pause

  \begin{block}{}
\# mkfs.ext4 /dev/md0 \\
\# cat /proc/mdstat
  \end{block}
  Esperamos a que sincronice.

  \begin{block}{}
\# mount /dev/md0 /mnt \\

\# echo bar $>$ /mnt/foo
  \end{block}

\end{frame}


%%%%%%%%%%%%%%%%%%%%%%%%%%%%%%%%%%%%%%%%%%%%%%%%%%%%%%%%%%%%%%%%%%%%%%%%%

\begin{frame}
  \frametitle{RAIDs: Algunas pruebas de funcionamiento del RAID (1)}

Reiniciamos el PC y vemos que está el RAID perfectamente funcional y lo podemos volver a montar:

\pause

  \begin{block}{}
\# cat /proc/mdstat \\
\# mount /dev/md0 /mnt \\
\# cat /mnt/foo
  \end{block}

\end{frame}


%%%%%%%%%%%%%%%%%%%%%%%%%%%%%%%%%%%%%%%%%%%%%%%%%%%%%%%%%%%%%%%%%%%%%%%%%

\begin{frame}
  \frametitle{RAIDs: Algunas pruebas de funcionamiento del RAID (y 2)}

Apagamos el PC y eliminamos uno de los dispositivos de bloques del RAID. Arrancamos de nuevo el PC.
Vemos que nos llega un mail informando de que el RAID está degradado:

  \begin{block}{}
\# mail
  \end{block}

\pause

Aún así, el RAID sigue estando operativo y se puede seguir utilizando:

  \begin{block}{}
\# cat /proc/mdstat \\
\# mount /dev/md0 /mnt \\

\# cat /mnt/foo \\
\# echo foo $>$ /mnt/bar
  \end{block}

\end{frame}


\subsection{LVMs}

%%%%%%%%%%%%%%%%%%%%%%%%%%%%%%%%%%%%%%%%%%%%%%%%%%%%%%%%%%%%%%%%%%%%%%%%%

\begin{frame}
  \frametitle{LVMs: Ejercicio 1}

\begin{itemize}
\item \textbf{Crear un dispositivo de bloques virtual más y configurar el raid y el nuevo dispositivo como dispositivos físicos.}
\end{itemize}

\pause

  \begin{block}{}
\# pvcreate /dev/md0 \\
\# pvcreate /dev/sdd
  \end{block}

{\footnotesize \alert{OJO:} Ahora somos los administradores los que tenemos que tener cuidado de \alert{no montar \texttt{/dev/md0}} ya que son dos tecnologías diferentes y no hay nada que nos lo impida.}

\end{frame}



\begin{frame}
  \frametitle{LVMs: Ejercicio 2}

\begin{itemize}
\item \textbf{Crear un grupo de volúmenes con los dos dispositivos físicos anteriores.}
\end{itemize}

\pause

  \begin{block}{}
\# vgcreate vg00 /dev/md0 /dev/sdd
  \end{block}

\end{frame}


\begin{frame}
  \frametitle{LVMs: Ejercicio 3}

\begin{itemize}
\item \textbf{Crear un volumen lógico de la mitad de tamaño del volumen}
\item \textbf{Formatear ese volumen lógico, montarlo en \texttt{/mnt} y crear algún fichero de prueba}
\end{itemize}

\pause

  \begin{block}{}
\# lvcreate -L1G -n foo vg00 \\
\# mkfs.ext3 /dev/vg00/foo \\
\# mount /dev/vg00/foo /mnt
  \end{block}

\end{frame}


\begin{frame}
  \frametitle{LVMs: Ejercicio 4}

\begin{itemize}
\item \textbf{Crear un snapshot del volumen lógico y montarlo en otro punto del sistema de ficheros}
\item \textbf{Comprobar que se pueden añadir y borrar ficheros en el snapshot sin modificar el volumen lógico.}
\item \textbf{¿Para qué puede ser útil esta funcionalidad?}
\end{itemize}

\pause

  \begin{block}{}
\# lvcreate -L512M -s -n foo-snap /dev/vg00/foo \\
\# mount /dev/vg00/foo-snap /mnt2
  \end{block}

\end{frame}

\begin{frame}
  \frametitle{LVMs: Ejercicio 4}

\begin{itemize}
\item \textbf{Formatear el volumen lógico con un sistema de ficheros que permita el crecimiento en caliente.}
\item \textbf{Aumentar el tamaño del volumen lógico y posteriormente el tamaño del sistema de ficheros}
\end{itemize}

\pause

  \begin{block}{}
\# lvremove /dev/vg00/foo-snap \\
\# lvremove /dev/vg00/foo \\
 
\# lvcreate -L512M -n foo vg00 \\

\# apt-get install xfsdump \\

\# mkfs.xfs /dev/vg00/foo \\
\# mount /dev/vg00/foo /mnt \\

\# lvresize -L+1G /dev/vg00/foo \\

\# xfs\_grow /mnt
  \end{block}

\end{frame}



\subsection{NFS}

\begin{frame}
  \frametitle{Ejercicio con NFS}

  \begin{itemize}
\item Añadir a la máquina virtual creada un nuevo interfaz de red de tipo `internal'
\item Crear una segunda máquina virtual con Debian y un interfaz de red de tipo `internal' conectado a la misma red que la anterior.
\item Instalar en una de ellas un servidor de NFS y servir un directorio que será montado en la otra máquina por la red interna.
  \end{itemize}

\end{frame}







%%%%%%%%%%%%%%%%%%%%%%%%%%%%%%%%%%%%%%%%%%%%%%%%%%%%%%%%%%%%%%%%%%%%%%%%%

% Portada final
\frame{
\maketitle
\begin{center}
\includegraphics[width=6cm]{format/gsyc-urjc}
\end{center}
}


%Fin de la presentacion
\end{document}


